% GENERATED FILE. Edit canonical lessons, then run npm run generate:books.
% canonical-source-hash: fnv1a64:d04c5315706d92a6
% canonical-lessons: JA-C517-ketsudan-suru, JA-C517-ketsui-suru, JA-C517-omoikiru, JA-C517-handan-suru, JA-C517-sentakushi

\chapter{Making Up One's Mind}
\label{ch:ja-making-up-one-s-mind}
\hlchaptermodality{517}

\begin{chapteropening}
\textbf{By the end of this chapter:} \emph{I can talk about deciding and hesitating.}

Complete the last lesson of chapter 517: I can talk about deciding and hesitating.
\end{chapteropening}

\section[\texorpdfstring{\ja{けつだんする} (ketsudan suru)}{けつだんする (ketsudan suru)}]{\ja{けつだんする} (ketsudan suru) — to make a firm decision}

\label{lesson:JA-C517-ketsudan-suru}



\begin{quote}
Before the new one: say the Japanese for the will to do something, then the Japanese for to tackle (a challenge).
\end{quote}

\subsection*{You'll want to know: \ja{けつだんする}}
\textbf{\ja{けつだんする}} — \emph{ketsudan suru} — \textquotedblleft{}to make a firm decision\textquotedblright{}.

\begin{quote}
\textbf{\ja{もし} \ja{はやく} \ja{けつだんしたら、きっぷが} \ja{やすいでしょう。}} — \emph{moshi hayaku ketsudanshitara, kippu ga yasui deshō} — \textquotedblleft{}If you decide early, the tickets will probably be cheaper.\textquotedblright{}
\end{quote}

Polite \textbf{\ja{けつだんします}}, \emph{ketsudanshimasu}; a \ja{する} verb. Past (polite): \textbf{\ja{しゃちょうが} \ja{けつだんしました。}}, \emph{shachō ga ketsudanshimashita}.

\begin{cousinweb}
\textbf{\ja{けつだんする}} (\emph{ketsudan suru}) is the Sino-Japanese (kango) noun \emph{ketsudan} with the native verb \emph{suru}, \textquotedblleft{}to do\textquotedblright{}: two Chinese characters read with their on'yomi, readings borrowed from Chinese. They are \emph{ketsu}, \textquotedblleft{}decide\textquotedblright{} (Middle Chinese \emph{kwet}) and \emph{dan}, \textquotedblleft{}severance\textquotedblright{} (\emph{dwanX}) (KANJIDIC2; Middle Chinese in Baxter's transcription).
\end{cousinweb}

\begin{grammarlens}[title={What you've built}]
One more word for saying what happens if.
\end{grammarlens}

\subsection*{Your turn}
\begin{itemize}
\raggedright
  \item \emph{Say it:} \emph{ketsudan suru}
  \item \emph{Say it:} \emph{ketsudan suru}, once more
  \item \emph{Say it:} \emph{moshi hayaku ketsudanshitara, kippu ga yasui deshō} — If you decide early, the tickets will probably be cheaper.
  \item \emph{Recall:} say the Japanese for the will to do something, then the Japanese for to tackle (a challenge), then say \emph{ketsudan suru} again
\end{itemize}

\begin{tcolorbox}[breakable,colback=teal!4,colframe=teal!35!black,title={Before you move on}]
What does \ja{けつだんする} mean? (\textquotedblleft{}To make a firm decision\textquotedblright{}.) Say it once more.
\end{tcolorbox}
\section[\texorpdfstring{\ja{けついする} (ketsui suru)}{けついする (ketsui suru)}]{\ja{けついする} (ketsui suru) — to resolve (to do)}

\label{lesson:JA-C517-ketsui-suru}



\begin{quote}
Before the new one: say the Japanese for to tackle (a challenge), then the Japanese for to make a firm decision.
\end{quote}

\subsection*{You'll want to know: \ja{けついする}}
\textbf{\ja{けついする}} — \emph{ketsui suru} — \textquotedblleft{}to resolve (to do)\textquotedblright{}.

\begin{quote}
\textbf{\ja{もし} \ja{けついしたら、さいごまで} \ja{がんばって} \ja{ください。}} — \emph{moshi ketsuishitara, saigo made ganbatte kudasai} — \textquotedblleft{}If you make up your mind, please keep going to the end.\textquotedblright{}
\end{quote}

Polite \textbf{\ja{けついします}}, \emph{ketsuishimasu}; a \ja{する} verb. Past (polite): \textbf{\ja{かれが} \ja{ひっこしを} \ja{けついしました。}}, \emph{kare ga hikkoshi o ketsuishimashita}.

\begin{cousinweb}
\textbf{\ja{けついする}} (\emph{ketsui suru}) is the Sino-Japanese (kango) noun \emph{ketsui} with the native verb \emph{suru}, \textquotedblleft{}to do\textquotedblright{}: two Chinese characters read with their on'yomi, readings borrowed from Chinese. They are \emph{ketsu}, \textquotedblleft{}decide\textquotedblright{} (Middle Chinese \emph{kwet}) and \emph{i}, \textquotedblleft{}idea\textquotedblright{} (\emph{'iH}) (KANJIDIC2; Middle Chinese in Baxter's transcription).
\end{cousinweb}

\begin{grammarlens}[title={What you've built}]
One more word for saying what happens if.
\end{grammarlens}

\subsection*{Your turn}
\begin{itemize}
\raggedright
  \item \emph{Say it:} \emph{ketsui suru}
  \item \emph{Say it:} \emph{ketsui suru}, once more
  \item \emph{Say it:} \emph{moshi ketsuishitara, saigo made ganbatte kudasai} — If you make up your mind, please keep going to the end.
  \item \emph{Recall:} say the Japanese for to tackle (a challenge), then the Japanese for to make a firm decision, then say \emph{ketsui suru} again
\end{itemize}

\begin{tcolorbox}[breakable,colback=teal!4,colframe=teal!35!black,title={Before you move on}]
What does \ja{けついする} mean? (\textquotedblleft{}To resolve (to do)\textquotedblright{}.) Say it once more.
\end{tcolorbox}
\section[\texorpdfstring{\ja{おもいきる} (omoikiru)}{おもいきる (omoikiru)}]{\ja{おもいきる} (omoikiru) — to give up all thought of}

\label{lesson:JA-C517-omoikiru}



\begin{quote}
Before the new one: say the Japanese for to make a firm decision, then the Japanese for to resolve (to do).
\end{quote}

\subsection*{You'll want to know: \ja{おもいきる}}
\textbf{\ja{おもいきる}} — \emph{omoikiru} — \textquotedblleft{}to give up all thought of\textquotedblright{}.

\begin{quote}
\textbf{\ja{もし} \ja{ねだんが} \ja{あがったら、くるまは} \ja{おもいきります。}} — \emph{moshi nedan ga agattara, kuruma wa omoikirimasu} — \textquotedblleft{}If the price goes up, I will give up on the car.\textquotedblright{}
\end{quote}

Polite \textbf{\ja{おもいきります}}, \emph{omoikirimasu}; a godan verb. Past (polite): \textbf{\ja{ちちが} \ja{くるまを} \ja{おもいきりました。}}, \emph{chichi ga kuruma o omoikirimashita}.

\begin{cousinweb}
\textbf{\ja{おもいきる}} (\emph{omoikiru}) is a native Japanese word (wago), inherited rather than borrowed from Chinese, built from \emph{omou}, \textquotedblleft{}to think\textquotedblright{}, and \emph{kiru}, \textquotedblleft{}to cut\textquotedblright{}; its kanji was attached later as a kun'yomi (JMdict, KANJIDIC2).
\end{cousinweb}

\begin{grammarlens}[title={What you've built}]
One more word for saying what happens if.
\end{grammarlens}

\subsection*{Your turn}
\begin{itemize}
\raggedright
  \item \emph{Say it:} \emph{omoikiru}
  \item \emph{Say it:} \emph{omoikiru}, once more
  \item \emph{Say it:} \emph{moshi nedan ga agattara, kuruma wa omoikirimasu} — If the price goes up, I will give up on the car.
  \item \emph{Recall:} say the Japanese for to make a firm decision, then the Japanese for to resolve (to do), then say \emph{omoikiru} again
\end{itemize}

\begin{tcolorbox}[breakable,colback=teal!4,colframe=teal!35!black,title={Before you move on}]
What does \ja{おもいきる} mean? (\textquotedblleft{}To give up all thought of\textquotedblright{}.) Say it once more.
\end{tcolorbox}
\section[\texorpdfstring{\ja{はんだんする} (handan suru)}{はんだんする (handan suru)}]{\ja{はんだんする} (handan suru) — to judge; to decide (after weighing it up)}

\label{lesson:JA-C517-handan-suru}



\begin{quote}
Before the new one: say the Japanese for to resolve (to do), then the Japanese for to give up all thought of.
\end{quote}

\subsection*{You'll want to know: \ja{はんだんする}}
\textbf{\ja{はんだんする}} — \emph{handan suru} — \textquotedblleft{}to judge; to decide (after weighing it up)\textquotedblright{}.

\begin{quote}
\textbf{\ja{もし} \ja{じかんが} \ja{あったら、よく} \ja{かんがえて} \ja{はんだんします。}} — \emph{moshi jikan ga attara, yoku kangaete handan shimasu} — \textquotedblleft{}If I have time, I will think it over and decide.\textquotedblright{}
\end{quote}

Polite \textbf{\ja{はんだんします}}, \emph{handanshimasu}; a \ja{する} verb. Past (polite): \textbf{\ja{いしゃが} \ja{はんだんしました。}}, \emph{isha ga handanshimashita}.

\begin{cousinweb}
\textbf{\ja{はんだんする}} (\emph{handan suru}) is the Sino-Japanese (kango) noun \emph{handan} with the native verb \emph{suru}, \textquotedblleft{}to do\textquotedblright{}: two Chinese characters read with their on'yomi, readings borrowed from Chinese. They are \emph{han}, \textquotedblleft{}judgement\textquotedblright{} (Middle Chinese \emph{phanH}) and \emph{dan}, \textquotedblleft{}judgement\textquotedblright{}, whose Middle Chinese reading for this sense does not match the voiced Japanese one, so none is cited (KANJIDIC2; Middle Chinese in Baxter's transcription).
\end{cousinweb}

\begin{grammarlens}[title={What you've built}]
One more word for saying what happens if.
\end{grammarlens}

\subsection*{Your turn}
\begin{itemize}
\raggedright
  \item \emph{Say it:} \emph{handan suru}
  \item \emph{Say it:} \emph{handan suru}, once more
  \item \emph{Say it:} \emph{moshi jikan ga attara, yoku kangaete handan shimasu} — If I have time, I will think it over and decide.
  \item \emph{Recall:} say the Japanese for to resolve (to do), then the Japanese for to give up all thought of, then say \emph{handan suru} again
\end{itemize}

\begin{tcolorbox}[breakable,colback=teal!4,colframe=teal!35!black,title={Before you move on}]
What does \ja{はんだんする} mean? (\textquotedblleft{}To judge; to decide (after weighing it up)\textquotedblright{}.) Say it once more.
\end{tcolorbox}
\section[\texorpdfstring{\ja{せんたくし} (sentakushi)}{せんたくし (sentakushi)}]{\ja{せんたくし} (sentakushi) — an option; a choice}

\label{lesson:JA-C517-sentakushi}



\begin{quote}
Before the new one: say the Japanese for to give up all thought of, then the Japanese for to judge.
\end{quote}

\subsection*{You'll want to know: \ja{せんたくし}}
\textbf{\ja{せんたくし}} — \emph{sentakushi} — \textquotedblleft{}an option; a choice\textquotedblright{}.

\begin{quote}
\textbf{\ja{もし} \ja{せんたくしが} \ja{ふたつ} \ja{あったら、どちらを} \ja{えらびますか。}} — \emph{moshi sentakushi ga futatsu attara, dochira o erabimasu ka} — \textquotedblleft{}If there were two options, which would you choose?\textquotedblright{}
\end{quote}

\begin{cousinweb}
\textbf{\ja{せんたくし}} (\emph{sentakushi}) is Sino-Japanese (kango): three Chinese characters read with their on'yomi, readings borrowed from Chinese. They are \emph{sen}, \textquotedblleft{}elect\textquotedblright{} (Middle Chinese \emph{sjwenX}), \emph{taku}, \textquotedblleft{}choose\textquotedblright{} and \emph{shi}, \textquotedblleft{}limb\textquotedblright{} (KANJIDIC2; Middle Chinese in Baxter's transcription).
\end{cousinweb}

\begin{grammarlens}[title={What you've built}]
One more word for saying what happens if.
\end{grammarlens}

\subsection*{Your turn}
\begin{itemize}
\raggedright
  \item \emph{Say it:} \emph{sentakushi}
  \item \emph{Say it:} \emph{sentakushi}, once more
  \item \emph{Say it:} \emph{moshi sentakushi ga futatsu attara, dochira o erabimasu ka} — If there were two options, which would you choose?
  \item \emph{Recall:} say the Japanese for to give up all thought of, then the Japanese for to judge, then say \emph{sentakushi} again
\end{itemize}

\begin{tcolorbox}[breakable,colback=teal!4,colframe=teal!35!black,title={Before you move on}]
What does \ja{せんたくし} mean? (\textquotedblleft{}An option; a choice\textquotedblright{}.) Say it once more.
\end{tcolorbox}
